\chapter{Centralização do alvo, perda de profundidade e estimação}
\label{cap:centralizacao}

Este capítulo trata de três camadas que ficam em volta da lei de controle e decidem o
comportamento quando a medida é ruim ou incompleta: o canal vertical (como manter o alvo
no centro da imagem com a câmera fixa), a aproximação sem profundidade (APPROACH) e a
qualidade da estimativa 3D.

\section{Canal vertical (\code{vertical\_mode})}
A câmera do x500\_depth é fixa (o gimbal opcional está na \autoref{sec:gimbal}). Centralizar o alvo exige dois canais: o yaw
(\code{yaw\_mode: face\_target}) cuida da horizontal, e o canal vertical decide a altitude.
A classe \code{control/vertical\_policy.cpp} substitui apenas a componente $z$ do comando de
\emph{qualquer} lei:
\begin{equation}
  v_z = v_{alvo,z} + k\,e_z,\qquad |v_z|\le v_{z,max},
\end{equation}
\begin{itemize}
  \item \code{los} (padrão): sem substituição. A lei segue o ponto de \emph{standoff} na
        linha de visada 3D e preserva o ângulo de elevação inicial. Um alvo 3 m acima
        continua acima.
  \item \code{match\_altitude}: $e_z=p_{rel,z}+h_{off}$, ou seja, mesma altitude do alvo
        mais um deslocamento.
  \item \code{center\_image}: $e_z=p_{rel,z}-d_h\tan\varepsilon$, com $d_h$ a distância
        horizontal e $\varepsilon$ a elevação do eixo óptico (montagem da câmera + arfagem medida,
        filtrada por um passa-baixas de 1ª ordem). Com câmera e corpo nivelados equivale a
        \code{match\_altitude}. Ao arfar para acelerar, a câmera aponta para baixo e o drone
        sobe para manter o alvo no centro.
\end{itemize}
Como é um laço P com \emph{feedforward} de velocidade, o erro de regime para alvo com
velocidade vertical constante é nulo, e a dinâmica do erro é de 1ª ordem com constante de
tempo $1/k$ (padrão $k=\SI{1}{s^{-1}}$), desde que a malha de velocidade da PX4 seja mais
rápida que isso.

\section{Câmera em gimbal}
\label{sec:gimbal}
Com \code{AERO\_CAMERA\_MOUNT=gimbal} o perseguidor é o modelo \code{x500\_gimbal\_depth}: a
mesma OakD-Lite (RGB + profundidade) sobre um gimbal de 3 eixos sob a frente do corpo, com a cadeia
\begin{equation}
  R_{BC} = R_z(\psi_g)\,R_x(\phi_g)\,R_y(-\theta_g),
\end{equation}
em que $\theta_g>0$ eleva o eixo óptico. O nó \code{gimbal\_controller} separa duas tarefas.
\begin{itemize}
  \item \textbf{Estabilização}: a cada passo (\SI{100}{Hz}) os ângulos das juntas resolvem
        $R_{BC}=R_{WB}^{\top}R_{WC}^{*}$, com $R_{WC}^{*}$ o eixo óptico desejado no mundo e o
        horizonte nivelado (eixo $y$ da câmera horizontal). A rotação do corpo é rejeitada sem
        atraso de filtro; sobra a malha de posição da junta no Gazebo (comando de velocidade,
        $K_p=\SI{40}{s^{-1}}$, $\tau\approx\SI{25}{ms}$).
  \item \textbf{Apontamento}: em \code{track} a direção desejada é a linha de visada
        $\hat p_{rel}-R_{WB}\,p_{pivô}$ (estimativa filtrada, corrigida pela posição do pivô) ou,
        sem profundidade, a LOS ENU da medida. Azimute e elevação dessa direção passam por um
        seguidor com \emph{inércia} (\code{inertia.enabled}, \textbf{ligado por padrão}):
        \begin{equation}
          v^{*}=\operatorname{sign}(e)\min\!\left(k|e|,\ \sqrt{2a_{max}|e|},\ \omega_{max}\right)+\dot\chi_{ff},
          \qquad |\dot v|\le a_{max},
        \end{equation}
        com $\omega_{max}=\SI{90}{\degree/s}$, $a_{max}=\SI{360}{\degree/s^2}$, $k=\SI{6}{s^{-1}}$ e
        $\dot\chi_{ff}$ a taxa da direção desejada filtrada ($\tau=\SI{0.2}{s}$), que zera o
        atraso de regime ao seguir um alvo em movimento uniforme. O termo $\sqrt{2a_{max}|e|}$ é o
        perfil de frenagem ótimo: não há sobressinal. Desligada, a direção desejada vai direto às
        juntas.
\end{itemize}
A atitude $R_{WB}$ tem de ser a mesma que o estimador usou. Com visão, a medida entra no ENU pela
atitude do EKF da PX4, então o gimbal também usa o EKF (o erro de rumo do EKF se cancela). No modo
omnisciente a posição relativa é a verdade do Gazebo, e o rumo do EKF chegou a diferir
\SI{8}{\degree} do verdadeiro no SITL; por isso \code{attitude\_source: auto} usa a atitude verdadeira
nesse modo. Antes da correção, o alvo ficava \SI{6}{\degree} fora do centro.

No SITL (alvo a \SI{5}{m}, \SI{1.5}{m/s} de lado, corpo com \code{yaw\_mode: hold}) o erro do eixo
óptico ficou abaixo de \SI{1}{\degree} no omnisciente e de \SI{1.1}{\degree} com visão, com o
gimbal girando \SI{61}{\degree} em yaw. Assim o laço do corpo só precisa manter a distância:
\code{yaw\_mode: hold}/\code{velocity} e \code{vertical\_mode: los}/\code{match\_altitude}. Com
\code{center\_image} a referência passa a ser o eixo neutro do gimbal (\code{neutral\_pitch\_deg},
TF \code{gimbal\_neutral}), o que mantém a junta de pitch perto do meio do curso; a compensação de
arfagem deve ser desligada, pois o gimbal já estabiliza. O modo \code{locked} trava as juntas e
reproduz a câmera fixa no mesmo modelo, o que permite comparar as duas montagens no mesmo voo.

\section{Aproximação sem profundidade (APPROACH)}
\label{sec:approach}
Quando o rastreador vê o alvo, mas a câmera de profundidade não fornece pontos válidos,
o supervisor entra em APPROACH. A câmera simulada (OakD-Lite) mede de \SIrange{0.2}{19.1}{m}; com
a ROI toda além do alcance (no Gazebo os pixels vêm $+\infty$), a linha de visada traz o limite
inferior $r_{min}=\SI{19.1}{m}$ e a distância usada abaixo é $\max(\hat r, r_{min})$. O APPROACH comanda velocidade ao longo da direção do alvo:
\begin{equation}
  \mathbf v = s\;\phi_{al}\;\phi\;\hat{\mathbf d},\qquad
  \phi_{al}=\operatorname{sat}_{[0,1]}\!\left(1-\frac{|\alpha|}{\alpha_{max}}\right),
\end{equation}
em que $\alpha$ é o desvio do alvo em relação ao eixo óptico. A primeira versão usava
$\phi=1$, isto é, avançava em velocidade de cruzeiro sem saber a distância, e causou
colisões quando a profundidade falhava com o alvo próximo. Hoje $\phi$ combina duas fontes.

\subsection{Distância por memória ou pelo tamanho da ROI}
Se existe medida 3D com até \code{range\_memory} segundos, usa-se a estimativa
conservadora $\hat r=\min(|\mathbf p|,|\mathbf p+\mathbf v\,\Delta t|)$. Sem ela, e
conhecido o tamanho físico $D$ do alvo, o tamanho angular $\theta$ da ROI dá
$\hat r=D/(2\tan(\theta/2))$. Com $\hat r$:
\begin{equation}
  \phi_r=\operatorname{sat}_{[-\rho,\,1]}\!\left(\frac{\hat r-(d+m)}{\Delta_{slow}}\right),
\end{equation}
que para em $d+m$ (\emph{standoff} + margem) e \textbf{recua} (até $\rho\,s$) mais perto que
isso. O recuo não depende do alinhamento, porque afastar-se é seguro com qualquer rumo.

\paragraph{De onde vem $D$.} A caixa do rastreador varia quadro a quadro e herda a escala da
caixa semeada na aquisição: medindo $D=\sqrt{wh}\,r/f$ com profundidade, o X4 deu de 0,32 a
\SI{0,55}{m}, conforme a execução. Um $D$ fixo erra a distância em $\approx25\,\%$. Por isso
$\theta$ é filtrado (\code{size\_filter}) e $D$ é \textbf{aprendido} enquanto há
profundidade: $D_k=2\,\hat r\tan(\theta/2)$, com média móvel de constante
\code{learn\_time}. O valor configurado (\code{target\_size}) só vale até a primeira
medida. Num teste com \code{target\_size}$=\SI{1,0}{m}$ (errado), o nó aprendeu
\SI{0,60}{m} (medição independente: \SI{0,59}{m}) e parou a \SI{5,9}{m} do alvo
(ponto de parada \SI{6}{m}).

\subsection{Expansão da ROI e tempo até o contato}
Mesmo sem distância, a \emph{evolução} do tamanho da ROI informa a taxa de aproximação.
Para um alvo de tamanho $D$ a uma distância $r$, $\theta\approx D/r$, logo
\begin{equation}
  \varepsilon\equiv\frac{d\ln\theta}{dt}=-\frac{\dot r}{r}=\frac1\tau,
\end{equation}
em que $\tau$ é o tempo até o contato \cite{lee2009ttc}, que não depende de $D$. O nó
estima $\varepsilon$ pela diferença de $\ln\theta$ entre linhas de visada consecutivas,
filtrada (\code{expansion\_filter}), e aplica
\begin{equation}
  \phi_\varepsilon=\begin{cases}
    \min\!\left(\beta,\ 1+g\,(-\varepsilon)\right), & \varepsilon\le0
       \ \text{(ROI estável ou encolhendo: alvo fugindo, acelera)}\\[2pt]
    \operatorname{sat}_{[0,1]}\!\left(\dfrac{\tau-\tau_{min}}{\tau_{slow}-\tau_{min}}\right),
       & \varepsilon>0\ \text{(ROI crescendo: freia e para em } \tau_{min})
  \end{cases}
\end{equation}
Combinação: $\phi=\min(\phi_r,\phi_\varepsilon)$ quando a distância conhecida já pede
desaceleração ($\phi_r<1$); caso contrário, $\phi=\phi_\varepsilon$. Perto do ponto de parada
o drone nunca acelera, mesmo com o alvo fugindo. Padrões: $\tau_{min}=\SI{1,5}{s}$,
$\tau_{slow}=\SI{6}{s}$, $\beta=1{,}5$, $g=\SI{5}{s}$.

\begin{atencao}
O ruído na caixa do rastreador entra em $\varepsilon$ por derivação. Com o filtro de
\SI{0,5}{s}, uma oscilação de 5\,\% no lado da ROI a \SI{10}{Hz} ainda produz
$|\varepsilon|$ da ordem de $0{,}1\,\si{s^{-1}}$ (equivalente a $\tau=\SI{10}{s}$). Esse é o
motivo para $\tau_{slow}$ não ser muito maior que alguns segundos.
\end{atencao}

\section{Qualidade da estimativa 3D}
Três problemas encontrados em voo, com suas causas:
\begin{enumerate}
  \item \textbf{Profundidade quase nunca pareada.} RGB e profundidade vêm de sensores
        distintos, sem sincronismo ($\approx$10 e \SI{13}{Hz} no SITL). O rastreador comparava
        a imagem só com o \emph{último} quadro de profundidade e exigia
        $|\Delta t|\le\SI{50}{ms}$, então descartava quase tudo (1545 linhas de visada contra
        17 medidas 3D numa sessão). Agora escolhe, num pequeno buffer, o quadro de carimbo
        mais próximo.
  \item \textbf{Profundidade perdida com inferência lenta.} Com OpenCL, o rastreador
        confirmava o alvo, mas 0\,\% das detecções tinham profundidade. A profundidade e a
        imagem estavam no mesmo grupo de \emph{callback}. Enquanto a rede processava um quadro
        ($\approx\SI{70}{ms}$), a fila de profundidade (2 quadros) só guardava os mais novos, e
        nenhum casava com o carimbo da imagem. A profundidade ganhou grupo próprio, e o OpenCL
        passou a ter 100\,\% de profundidade. O erro \code{CL\_INVALID\_WORK\_GROUP\_SIZE} do
        \code{gemm} é inofensivo: vem do núcleo do OpenCV, que cai para a CPU com resultado
        correto (a rede ViT dá a mesma saída na GPU e na CPU, com diferença de $10^{-6}$).
  \item \textbf{Ponto ``tremendo'' no modo omnisciente.} A posição relativa era obtida no
        referencial do corpo (atitude \emph{real} do Gazebo) e depois girada para ENU com a
        atitude \emph{estimada} pelo EKF, num instante diferente. O erro resultante,
        $\approx r\,\delta\psi$, chegou a \SI{0,19}{m} em média e \SI{0,38}{m} no máximo. Agora a
        diferença $\mathbf p_{alvo}-\mathbf p_{perseguidor}$ é tirada direto no referencial
        do mundo, com as duas poses no mesmo instante. O filtro de Kalman (ajustado para a
        visão) ainda atrasava $\approx\SI{7}{cm}$ nas reversões do alvo. Como o omnisciente é a
        verdade do simulador, ele agora \textbf{não filtra}: publica a pose exata e a
        velocidade pela diferença de poses consecutivas (erro medido: \SI{0}{mm}).
\end{enumerate}
